\chapter{逆向工程}
\label{chap:reverse-engineering}
\label{ch:reverse}

\nt{本章的主线是\textbf{实物 $\to$ 测量数据 $\to$ 点云/网格 $\to$ CAD模型}。复习时要能说明每一步解决什么问题、输出什么数据，以及误差如何传到下一步。扫描得到的密集点云并不自动具有可编辑的设计特征。}

\section{从实物到数字模型}
\dfn{逆向工程}{逆向工程（Reverse Engineering）以已有实物为起点，通过测量、数据处理和几何重建获得数字模型，为检验、复制、修改及再设计提供依据。正向工程从设计模型制造实物，逆向工程从实物恢复模型；两者可以在产品开发过程中反复衔接。}

基本流程为：数据获取、预处理、多边形化、分割与曲面拟合、CAD模型构建。不同目标不一定走完全部步骤：尺寸检验可将点云与设计模型配准后比较；3D打印可使用修复后的水密三角网格；需要尺寸驱动修改、布尔运算或工艺特征识别时，通常还要恢复解析曲面、自由曲面及其拓扑关系。
\reviewfig[.92]{ch06-workflow.pdf}{逆向工程的数据流与不同用途。摘自第6章讲义，PDF第5页。}{fig:reverse-workflow}
\sourcepage{第6章}{3--5}

\section{数据采集：先判断测什么，再选择测量方法}
\subsection{接触式、非接触式与内部结构获取}
\begin{table}[H]
\centering\small
\begin{tabularx}{\linewidth}{@{}p{2.4cm} X X@{}}
\toprule
方法 & 原理与优势 & 限制与选用依据 \\
\midrule
三坐标测量机（CMM） & 探头随正交直线轴移动，获取离散位置；适于精密尺寸和几何要素检测。 & 逐点采样较慢；需考虑探头可达性、测力和探针半径补偿。 \\
关节臂式测量 & 由关节角及机构参数求末端探头位置，移动方便、姿态灵活。 & 误差随结构、标定和姿态变化；接触可能使柔软表面变形。 \\
光学表面扫描 & 利用反射光和成像几何快速获得密集点云。 & 受遮挡、透明、反光、吸光和边缘混合像素影响，往往需多视角扫描。 \\
体数据获取 & CT等获取内部切片，分割后重建内外几何；MRI以磁共振信号成像。 & 图像对比度、分辨率和分割误差都会影响几何；应按材料及目标结构选用。 \\
破坏性断层扫描 & 逐层去除材料并拍摄截面，再重建三维形状。 & 能暴露内部结构，但消耗原件，不能用于必须保留的样品。 \\
\bottomrule
\end{tabularx}
\end{table}

接触式测量常对表面颜色不敏感，但不是所有孔槽都能测到：探头必须能够进入且与目标表面接触。若仪器记录的是球形探针\textbf{中心}，还需沿接触法向补偿球半径，才能得到表面点。非接触测量能避免探针压变形，却仍有视线与表面光学性质的限制。

主动光学测量主动投射点、线或结构光图案；被动视觉利用已有光照下的图像。两者都可能运用三角几何。讲义将CT与MRI放在内部形状获取模块中，\textbf{不能据此把MRI理解成X射线透射成像}；两者物理机理不同。本章只比较其几何数据获取用途。
\sourcepage{第6章}{6--10、16--23}

\subsection{三角测距：由光路几何求深度}
已知光源与相机的相对位置，测得光斑在图像上的位置后，利用视线与投射光线的交会求三维点。图\ref{fig:reverse-triangulation}中，$b$ 为基线长度，$f$ 为焦距，$p_x$ 为图像平面的有向距离，$\theta$ 为投射角；按图中的符号约定，有
\[
\frac{z}{f}=\frac{x}{p_x},\qquad
\tan\theta=\frac{b-x}{z}.
\]
将 $x=zp_x/f$ 代入第二式，得到
\[
\boxed{z=\frac{fb}{p_x+f\tan\theta}},\qquad
x=b-z\tan\theta.
\]
这里 $p_x$ 是\textbf{物理长度}，若读数以像素表示，须用像元尺寸或标定后的内参统一单位。不同图像坐标正方向会改变式中符号，应从光路关系推导，不能脱离坐标约定套用。
\reviewfig[.68]{ch06-triangulation.pdf}{主动三角测距的成像几何。摘自第6章讲义，PDF第11页。}{fig:reverse-triangulation}

投射一条线光束可以一次获得一条截线；扫描该线可覆盖更大区域。结构光将空间编码图案投向物体，通过图案对应关系建立相机视线与投射方向的联系。双目视觉则从两幅图像寻找同名点，以两条视线求三维位置；\textbf{对应点匹配和几何标定}是关键。

\ex{三角测距的单位检查}{按图中约定，若 $f=20\,\mathrm{mm}$、$b=100\,\mathrm{mm}$、$\theta=45^\circ$，测得 $p_x=5\,\mathrm{mm}$，则
\[
z=\frac{20\times100}{5+20\times1}=80\,\mathrm{mm},\qquad
x=100-80=20\,\mathrm{mm}.
\]
检查 $z/f=x/p_x=4$，两条几何关系一致。本例说明方法，不代表特定扫描仪的参数。}
\sourcepage{第6章}{11--14}

\subsection{时差测距与多视角采集}
时差法记录光信号发射后往返目标的时间差 $\Delta t$，在传播速度近似为 $c$ 时，
\[
\boxed{D\approx\frac{c\Delta t}{2}}.
\]
除以2是因为计时覆盖\textbf{往返}距离。该原理适于远距离和较大范围的扫描；实际还需扣除系统延迟，并考虑测时误差及介质传播速度。

一次表面扫描通常只能得到朝向传感器且可见的部分。凹槽、背面和遮挡处需要改变视角；多组数据必须先配准到共同坐标系再融合。边缘处的反射与遮挡容易产生异常值，增加扫描视角并不能自动消除这些误差。
\sourcepage{第6章}{15--16}

\section{点云预处理：邻域、法向、异常点与光顺}
\subsection{邻域决定局部几何的尺度}
\begin{itemize}
\item \textbf{$k$近邻：}选择距离目标点最近的 $k$ 个点，控制样本数；点密度变化时，邻域的物理大小也会变化。
\item \textbf{$\varepsilon$球邻域：}选择到目标点距离不超过 $\varepsilon$ 的点，控制物理尺度；稀疏区可能没有足够点。
\item \textbf{网格环邻域：}由边的连接关系定义。1环是与该顶点直接相连的邻点；更高环由拓扑距离扩展，需说明是否包含内层邻域。
\end{itemize}
\reviewfig[.48]{ch06-neighborhoods.pdf}{点云的距离邻域与网格的环邻域。摘自第6章讲义，PDF第24页。}{fig:reverse-neighborhoods}
邻域过小，估计容易受噪声影响；邻域过大，会跨越棱边、混合不同曲面。尺度应结合采样间距、曲率和最小保留特征选取。
\sourcepage{第6章}{24}

\subsection{用局部平面拟合法向：为什么取最小特征值}
设邻域点为列向量 $\mathbf q_i$，非负权重为 $w_i$，且 $\sum_iw_i>0$。以单位法向 $\mathbf n$ 和常数 $d$ 表示平面 $\mathbf n^{\mathsf T}\mathbf q=d$，最小化点到平面的平方距离：
\[
\min_{\mathbf n,d}\sum_iw_i(\mathbf n^{\mathsf T}\mathbf q_i-d)^2,
\qquad \mathbf n^{\mathsf T}\mathbf n=1.
\]
对 $d$ 求导得
\[
d=\mathbf n^{\mathsf T}\bar{\mathbf q},\qquad
\bar{\mathbf q}=\frac{\sum_iw_i\mathbf q_i}{\sum_iw_i}.
\]
即最优平面经过加权质心。代回目标函数，定义
\[
\mathbf C=\sum_iw_i(\mathbf q_i-\bar{\mathbf q})
(\mathbf q_i-\bar{\mathbf q})^{\mathsf T},
\]
目标变为 $\min\mathbf n^{\mathsf T}\mathbf C\mathbf n$。用拉格朗日乘子处理单位长度约束，有
\[
\boxed{\mathbf C\mathbf n=\lambda\mathbf n}.
\]
$\mathbf C$ 对称半正定，故\textbf{最小特征值对应的单位特征向量}给出法向：点集沿这个方向的离散程度最小，最接近薄平面的厚度方向。协方差是否除以权重和，不影响特征向量。

\nt{平面拟合只确定 $\mathbf n$ 与 $-\mathbf n$ 这一对方向，法向朝外还是朝内需要额外统一。如果邻域几乎共线，最小特征方向不唯一；若跨越两个面，拟合平面会偏离两侧真实法向。}

网格三角面的法向可由两条边叉乘并单位化得到。顶点法向可对相邻面的法向作加权平均后\textbf{再次单位化}，常用面积或角度权重。曲率可由邻域二次曲面拟合或离散微分几何估计；曲率比法向更依赖采样与噪声。
\sourcepage{第6章}{25--30}

\subsection{异常点剔除与噪声光顺的区别}
异常点是明显不符合局部表面形状或采样分布的观测；噪声则可能分布在所有有效点上。先删除离群点，再在需要时光顺，可避免异常点把局部平均拉偏。

讲义给出局部平面距离和邻域包围球两类判断。其共同思想是：把待判断点的偏差与邻域自身的尺度比较，而不是使用与单位无关的固定数值。例如，设待判断点到局部平面的距离为 $d_p$，其余邻点的平均残差为 $\bar d$，可以比较 $d_p/(d_p+\bar d)$ 与阈值；分母为零时须单独处理。锐边点不服从单一平面，\textbf{不能只凭平面残差大就删除}。

拉普拉斯光顺把点移向邻域加权中心。设 $\sum_{j\in N(i)}w_{ij}=1$，则
\[
\Delta\mathbf p_i=\sum_{j\in N(i)}w_{ij}
(\mathbf p_j-\mathbf p_i),\qquad
\boxed{\mathbf p_i^{\mathrm{new}}=\mathbf p_i+\lambda\Delta\mathbf p_i}.
\]
均匀权重取 $w_{ij}=1/|N(i)|$。在 $0\leq\lambda\leq1$ 且权重非负时，新点位于原点与邻域中心之间。迭代应使用同一轮旧坐标计算所有新点，否则结果会依赖更新顺序。

\ex{一次光顺更新}{若某点 $\mathbf p=(0,0,0.1)$，四个邻点均在 $z=0$ 上，且邻点坐标平均为 $(0,0,0)$，取均匀权重及 $\lambda=0.5$，则 $\Delta\mathbf p=(0,0,-0.1)$，更新为 $(0,0,0.05)$。噪声高度减半，但真实的小凸起也可能同样被削弱。}

\nt{拉普拉斯光顺可能使物体\textbf{几何收缩}、体积减小并模糊棱边；它通常不删除顶点，不能把“收缩”理解为点数减少。减少点数属于数据缩减。保特征光顺需限制跨棱边平均、约束边界或采用其他权重与模型。}
\sourcepage{第6章}{31--34}

\section{配准与融合：寻找同一表面的共同坐标}
\subsection{有对应点时的刚体最小二乘配准}
本节统一采用与第2章一致的\textbf{列向量}形式，目标点与源点满足 $\mathbf q_i\approx\mathbf R\mathbf p_i+\mathbf t$。遇到原讲义中的行向量右乘形式，应将齐次矩阵整体转置，并相应改变复合变换的乘法次序。

刚体配准不允许缩放或变形，求解
\[
\min_{\mathbf R,\mathbf t}\sum_{i=1}^n
\|\mathbf R\mathbf p_i+\mathbf t-\mathbf q_i\|^2,
\qquad
\mathbf R^{\mathsf T}\mathbf R=\mathbf I,\quad \det\mathbf R=1.
\]
理想无噪声下，至少三对对应且不共线的点可确定空间刚体变换。先求质心 $\bar{\mathbf p},\bar{\mathbf q}$，并去中心化：
\[
\mathbf p_i'=\mathbf p_i-\bar{\mathbf p},\quad
\mathbf q_i'=\mathbf q_i-\bar{\mathbf q},\quad
\boxed{\mathbf t=\bar{\mathbf q}-\mathbf R\bar{\mathbf p}}.
\]
构造交叉协方差矩阵并作奇异值分解：
\[
\mathbf H=\sum_i\mathbf p_i'(\mathbf q_i')^{\mathsf T}
=\mathbf U\boldsymbol\Sigma\mathbf V^{\mathsf T}.
\]
正确的旋转为
\[
\boxed{\mathbf R=\mathbf V\,
\operatorname{diag}(1,1,\det(\mathbf V\mathbf U^{\mathsf T}))
\mathbf U^{\mathsf T}}.
\]
当 $\det(\mathbf V\mathbf U^{\mathsf T})=1$ 时，简化为讲义中的 $\mathbf V\mathbf U^{\mathsf T}$。行列式修正防止把镜像反射误当作旋转；只有正交约束还不够。
\sourcepage{第6章}{35--40}

\subsection{ICP：对应关系与变换交替更新}
\dfn{迭代最近点法}{ICP（Iterative Closest Point）在已有初始对齐的基础上，交替寻找两组数据的对应点并估计刚体变换，以逐步减小重叠区域的配准误差。最近点是算法估计的对应点，不一定是真实的同一物理点。}
\begin{enumerate}
\item \textbf{初始配准：}借助标记点、人工特征对应或其他粗配准，使两组数据大致重叠。
\item \textbf{寻找点对：}将源点按当前变换映射到目标系，对每个源点寻找目标集最近点。
\item \textbf{拒绝错误点对：}剔除距离过大、法向差异过大或不属于重叠区域的对应。讲义的双向检查将 $\mathbf p_i$ 匹配到 $\mathbf q_i$，再由 $\mathbf q_i$ 找回源集最近点 $\widehat{\mathbf p}_i$；若 $\|\mathbf p_i-\widehat{\mathbf p}_i\|$ 超过阈值则拒绝。
\item \textbf{更新变换：}利用保留点对求最小二乘刚体变换，并与当前变换复合。
\item \textbf{判断终止：}误差下降量、变换增量足够小，或达到迭代次数上限时停止。
\end{enumerate}
\reviewfig[.73]{ch06-icp.pdf}{ICP的对应点及双向最近点检查。摘自第6章讲义，PDF第41页。}{fig:reverse-icp}

ICP是局部优化，初值偏差大、重叠少、存在离群点或几何退化时可能失败。例如一大片无纹理平面能很好限制法向位移，却难以可靠约束面内滑动；周期结构还可能配到错误的重复单元。因此“迭代收敛”只说明当前目标变化很小，\textbf{不能证明对齐位置正确}。

配准确定坐标关系，\textbf{融合}才进一步处理重叠数据的重复采样、噪声与连接。质量检验还应使用规定的基准对齐：任意最佳拟合可能把一部分制造偏差分摊到整个零件上。
\sourcepage{第6章}{40--42、72}

\subsection{数据缩减：用较少数据保留必要几何}
均匀体素抽样把空间划分为网格单元，以每个单元中的一个代表点或质心代替所有点，速度快，但可能损失薄壁、边界和小圆角。自适应缩减利用法向变化、局部拟合误差或曲率控制采样：平坦区域可少保留点，高曲率与特征处应较密。八叉树或BSP可组织这种分层空间划分。

\textbf{缩减的约束是允许几何误差，不是只追求压缩比例。}选取代表点后，质心未必仍位于原始曲面上；缩减结果还要检查边界与关键尺寸。
\sourcepage{第6章}{43--45}

\section{网格重建与模型修复}
\subsection{点坐标之外，还需要邻接关系}
三角网格用顶点、边和三角面表达表面的分段线性近似。点云只有采样位置，网格还提供哪些点相邻、哪些面共享边等拓扑信息。半边结构将一条边分成两个有向半边，以后继半边、反向半边、所属面和端点等链接支持局部遍历。

\dfn{流形网格与水密网格}{二维流形曲面在内部点附近像一个圆盘，在边界点附近像一个半圆盘。内部边通常恰有两个相邻三角面，边界边只有一个；顶点邻域还需形成单一的环或扇。水密表面没有开放边界，常用于描述实体外壳；作为有效实体边界，还需检查自相交、方向一致性等条件。}
\reviewfig[.83]{ch06-nonmanifold.pdf}{非流形顶点与非流形边。仅检查每条边的相邻面数，不能发现全部顶点问题。摘自第6章讲义，PDF第47页。}{fig:reverse-nonmanifold}

\nt{有边界的曲面也可以是流形，因此\textbf{流形不等于水密}。另一方面，每条边恰有两个入射面，也不保证顶点邻域正确，更不保证嵌入空间中没有自相交。}
\sourcepage{第6章}{46--49}

\subsection{三类重建思路}
\begin{enumerate}
\item \textbf{增量表面生长：}从种子三角形或种子边出发，沿当前边界不断添加三角形；关键是选择正确的邻点，避免跨越间隙或连接到错误表面。
\item \textbf{隐式曲面提取：}先用有符号距离、移动最小二乘（MLS）或径向基函数（RBF）等描述隐式曲面，再用等值面提取方法生成网格。讲义所列Marching Cubes通常称为移动立方体法。
\item \textbf{基于体积的方法：}对点集构造三维Delaunay四面体剖分，再依据内外判断保留边界。表面网格的三角形与体剖分的四面体不能混淆。
\end{enumerate}
\sourcepage{第6章}{50--52}

修复通常包括合并重复补丁、统一面方向、消除非流形单元与自相交、补洞、光顺和简化。讲义的补洞流程是：先对洞口作粗三角剖分，再细化以匹配周围采样密度，最后光顺使其与邻域几何衔接。\textbf{补洞恢复的是推测的表面}，缺测区域过大时应优先补测；若孔本来就是设计中的通孔，不能把真实孔口当成缺陷封住。
\sourcepage{第6章}{53--55}

\section{分割、曲面拟合与CAD模型恢复}
\subsection{为什么不对整个点云直接拟合一个曲面}
机械零件常由平面、圆柱面、圆锥面、球面、拉伸面、旋转面与局部过渡面构成。把所有数据强行拟合为一张自由曲面，会掩盖原来的几何结构，并使尺寸修改和工艺处理困难。\textbf{分割}将数据划成分别对应自然曲面的区域，为曲面类型识别与拟合建立依据。

基于边缘的分割利用法向突变、曲率峰值或曲率跳跃寻找边界；区域生长从种子区域出发，逐步加入符合当前曲面模型或相似几何性质的相邻点。区域扩大后可能需要重新识别曲面类型：大半径球面的一小块看起来近似平面，不能仅凭很小的局部区域就确定全局类型。边缘法、区域法和人工编辑可以结合。
\sourcepage{第6章}{56--61}

\subsection{工程曲面、自由曲面和设计约束}
\begin{itemize}
\item \textbf{工程对象：}优先拟合主要解析面或规则扫掠面，检查残差，再延伸、求交、修剪得到边界，最后补回圆角等次要过渡面。简单模型满足误差要求时，不必用更复杂的自由曲面代替。
\item \textbf{自由形状对象：}可先建立曲线网络，将区域划分为若干曲面片，再拟合B样条曲面，并处理片间连续性；自动划分后仍可能需要人工调整。
\item \textbf{联合约束：}平行、垂直、同轴、同心、对称及切线连续等应在拟合或后处理时显式考虑。分别拟合每个面，即使各自残差很小，也未必保持原设计关系。
\end{itemize}

残差小只说明模型接近\textbf{测量数据}，不自动说明它恢复了原始设计。需要区分实物上的制造误差、测量噪声和设计意图；“按实物复制”与“恢复理想尺寸”会得到不同模型。
\sourcepage{第6章}{62--67}

\ex{带孔支架的逆向建模方案}{目标是获得可修改孔径、孔距的支架CAD模型。可先扫描可见面，并对孔轴、基准面等关键要素进行补充测量；去除异常点、粗配准和ICP后，融合数据并检查缺测处。随后分割底面、侧面与圆柱孔壁，拟合平面和圆柱面，施加平行、垂直与同轴等有依据的约束，再求交修剪、恢复圆角，缝合为实体。最后对照点云检查整体偏差，并单独检查孔径和孔距。仅输出STL虽然能表示外观，却不能直接提供这些参数化设计特征。}

\Needspace{11cm}
\section{复习自测与易错点}
\begin{enumerate}
\item 点云、三角网格和B-rep模型依次多了什么信息？\textit{答：网格增加连接与邻接关系；B-rep进一步用几何曲面、边界曲线及拓扑关系表达模型。}
\item 局部PCA为何用最小特征值对应方向作法向？\textit{答：沿该方向的方差最小，即点到局部平面的平方残差最小。}
\item 拉普拉斯光顺后顶点数是否必然减少？\textit{答：否。普通光顺改变坐标，缩减才改变数据量。}
\item 三对共线对应点能否唯一确定空间刚体旋转？\textit{答：不能，绕该直线的旋转不能由这些对应点唯一确定。}
\item 为什么SVD配准需要检查行列式？\textit{答：正交矩阵也可能包含镜像反射，刚体旋转须满足 $\det\mathbf R=1$。}
\item ICP收敛是否表明配准正确？\textit{答：否。它可能停在局部错误匹配，需检查初值、重叠、残差分布及基准。}
\item 网格边只连接一个面，一定非流形吗？\textit{答：否。它可能是合法流形曲面的边界边。}
\item 自动补洞能否作为未知表面真实形状的证据？\textit{答：不能。补洞属于依据周围数据的几何推测。}
\end{enumerate}
\nt{按“\textbf{采集方式和盲区 $\to$ 误差与预处理 $\to$ 配准与融合 $\to$ 网格修复 $\to$ 曲面及设计约束}”解释完整案例。区分测量精度、采样密度和重建误差；点更多或残差更小，都不必然意味着恢复的设计更正确。}
